\begin{pastexam}[label={px:ncku-111-prog-01}]{成大資訊聯招 111 學年度 第 1 題}
（程式填空）

\textbf{[4\%]} The following C program converts a preorder expression to its postorder form. For
example, "+AB++ABC" is converted to "AB+AB+C+" with the program. Please fill in the missing parts of
the program in C language.
\begin{minted}[fontsize=\small]{c}
#include<stdio.h>
#include<stdlib.h>

#define MAX_INPUT 100000

typedef struct _Stack_node{
    char val;
    struct _Stack_node *prev;
} Stack_node;

typedef struct _Stack{
    Stack_node *top;
} Stack;

Stack *new_stack(){
    Stack *stk = malloc(sizeof(Stack));
    stk->top = NULL;
    return stk;
}

char Stack_pop(Stack *stk){
    char value = stk->top->val;
    Stack_node *tmp = stk->top;
    /** TODO (A) **/
    free(tmp);
    return value;
}

void Stack_push(Stack *stk, char value){
    Stack_node *new_node = malloc(sizeof(Stack_node));
    new_node->val = value;
    /** TODO (B) **/
    stk->top = new_node;
    return;
}

int Stack_is_empty(Stack *stk){
    return stk->top == NULL;
}

void pre_to_post(char *);
void post_to_pre(char *);

int main(){
    char input[MAX_INPUT];
    scanf("%s", input);
    pre_to_post(input);
    return 0;
}

void pre_to_post(char *input){
    Stack *stk = new_stack();
    for (int i = 0; input[i] != '\0'; ++i){
        switch(input[i]){
            case '+':
            case '-':
            case '*':
            case '/':
            case '^':
            case '%':
                Stack_push(stk, input[i]);
                break;
            default:
                printf("%c", input[i]);
                while (!Stack_is_empty(stk) && stk->top->val == 'A'){
                    char top = Stack_pop(stk);
                    printf("%c", Stack_pop(stk));
                }
                if (!Stack_is_empty(stk)) Stack_push(stk, 'A');
        }
    }
    printf("\n");
}
\end{minted}
\end{pastexam}

\begin{pastsoln}
答案：(A) \texttt{stk->top = tmp->prev;}（或 \texttt{stk->top = stk->top->prev;}）；
(B) \texttt{new\_node->prev = stk->top;}。

\textbf{堆疊的結構}。 這是用鏈結串列做的堆疊（資結書第~6.1~節）：\texttt{top} 指向頂端節點，每個節點的
\texttt{prev} 指向它下面那一個。
\begin{itemize}
\item \textbf{(A) pop}：先記下頂端的值和節點 \texttt{tmp}，在 \texttt{free(tmp)} 之前要把
\texttt{top} 移到下一個節點，否則 \texttt{top} 會指向已經釋放的記憶體。 所以填
\texttt{stk->top = tmp->prev;}。 順序不能和 \texttt{free} 對調：\texttt{free(tmp)} 以後就不能再讀
\texttt{tmp->prev}。
\item \textbf{(B) push}：新節點要壓在原本的頂端上面，所以它的 \texttt{prev} 指向原本的頂端，
\texttt{new\_node->prev = stk->top;}，下一行再把 \texttt{top} 改成新節點。 這兩行順序也不能反，先改
\texttt{top} 就找不回原本的頂端。 堆疊是空的時候 \texttt{stk->top} 是 \texttt{NULL}，新節點的
\texttt{prev} 也就是 \texttt{NULL}，不用另外處理。
\end{itemize}

\textbf{這個轉換在做什麼}。 前序式的每個運算子後面接著它的兩個運算元。 程式把運算子推進堆疊；字元
\texttt{'A'} 在堆疊裡是一個\textbf{記號}，表示「它下面那個運算子的左運算元已經輸出了」（運算元本身
不會進堆疊，所以這個記號不會和輸入裡的運算元 A 混淆）。 讀到一個運算元：
\begin{enumerate}
\item 先印出它。
\item 只要頂端是記號，就表示某個運算子的右運算元剛剛完成，彈出記號、再彈出那個運算子並印出（後序裡
運算子在兩個運算元之後）。 印出運算子代表一個更大的運算元也完成了，所以繼續檢查。
\item 停下來時，若堆疊不是空的，頂端是某個運算子，它的左運算元剛完成，推一個記號。
\end{enumerate}

\textbf{追蹤 \texttt{+AB++ABC}}（堆疊由底到頂，$\triangle$ 代表記號 \texttt{'A'}）：
\begin{center}
\small
\begin{tabular}{clll}
\toprule
讀入 & 動作 & 堆疊 & 已輸出 \\
\midrule
\texttt{+} & 推入 & \texttt{+} & \\
\texttt{A} & 印 A；頂端不是記號；推記號 & \texttt{+}$\triangle$ & \texttt{A} \\
\texttt{B} & 印 B；彈出記號和 \texttt{+}，印 \texttt{+}；堆疊空 & （空） & \texttt{AB+} \\
\texttt{+} & 推入 & \texttt{+} & \texttt{AB+} \\
\texttt{+} & 推入 & \texttt{++} & \texttt{AB+} \\
\texttt{A} & 印 A；推記號 & \texttt{++}$\triangle$ & \texttt{AB+A} \\
\texttt{B} & 印 B；彈出記號和 \texttt{+}，印 \texttt{+}；推記號 & \texttt{+}$\triangle$ & \texttt{AB+AB+} \\
\texttt{C} & 印 C；彈出記號和 \texttt{+}，印 \texttt{+}；堆疊空 & （空） & \texttt{AB+AB+C+} \\
\bottomrule
\end{tabular}
\end{center}
輸出 \texttt{AB+AB+C+}，和題目的例子一致。 注意題目的輸入 \texttt{+AB++ABC} 其實是兩個運算式
（\texttt{+AB} 和 \texttt{++ABC}）接在一起，程式照樣逐段轉換。 單一運算式 \texttt{*+AB-CD} 會轉成
\texttt{AB+CD-*}，和資結書第~6.3.5~節的前序轉後序結果相同。

另外，\texttt{while} 迴圈裡的 \texttt{char top = Stack\_pop(stk);} 彈出的是記號，變數 \texttt{top} 沒有
再用到；\texttt{post\_to\_pre} 只有宣告沒有定義，也沒有被呼叫，編譯和執行都沒有影響。
\end{pastsoln}
